\newpage
\section{Measurement Framework}
\label{sec:method_protocol}


\subsection{Objectives and Scope}
\label{subsec:measurement-objective}


% --- Block A: Measurement Objective ---

The conditional controller proposed in this thesis assumes that optional RAG modules
consume additional energy per request. The CITI KMS system described in
previously has no existing energy benchmarks nor toggled modules. Before designing a
controller, we must first quantify the energy cost of module activation.

To produce this evidence, the measurement framework records estimated energy consumption on each
of the three infrastructure nodes that compose the RAG pipeline:

\begin{figure}[h!]
    \centering
        \centering
        \includegraphics[width=0.8\textwidth]{energy-measurement.png}
        \caption{Node-level energy measurement architecture}
        \label{fig:energy-measurement}
\end{figure}

\begin{itemize}
    \item \textbf{embedding node}, which encodes queries and when reranking is
          active, runs a cross-encoder model;
    \item \textbf{vector-database node} (Milvus), which executes dense or hybrid
          similarity search;
    \item \textbf{LLM node}, which generates the final answer and when
          HyDE is active, produces a hypothetical document.
\end{itemize}

On each node, CPU and GPU energy are read from hardware counters at the start
and end of every request. The reported quantity is total energy per request in
joules, summed across all three nodes.

% writing_note : we need to explain why only CPU and GPU --> if not mentionned in chapter 2
% SOTA that GPU and CPU account for most energy usage

The primary comparison is between C0 (minimal pipeline: all modules off, dense search)
and C7 (the current CITI KMS configuration: all modules on, hybrid search). To enable
finer-grained attribution and provide training data for the controller, we extended the
comparison to all eight possible configurations (C0--C7) defined in
Section~\ref{subsec:baseline-config}. These form a complete $2^3$ factorial design over
three binary module axes: HyDE, reranking, and search type.

\begin{figure}[h!]
    \centering
        \centering
        \includegraphics[width=0.8\textwidth]{citi_config_tree_wide.png}
        \caption{$2^3$ factorial configuration space}
        \label{fig:config-tree}
\end{figure}

Energy alone does not suffice to evaluate module activation. A configuration that saves
energy but doubles response time or degrades answer correctness is not useful. Therefore,
each session also records latency and quality metrics.

\subsection{Unit of Analysis}
\label{subsec:unit-of-analysis}

The unit of analysis is a \emph{session}: one query executed under one fixed configuration
until an answer is produced. Every measurement (energy, time to first token, total
time, token counts, quality scores, etc.) is recorded under a unique
\texttt{session\_id}. A session is identified by the tuple \texttt{(query\_id, config\_id)}.

Sessions are executed sequentially with no concurrent load, independently of
configuration order. Each measurement therefore reflects the complete cost of
answering one query under one configuration.

% writing_note : keep this commented - do not delete.
% In total, the experimental dataset comprises 750~queries executed under all
% 8~configurations, yielding 6{'}000~sessions.

In total, each session records approximately 47~fields spanning energy, timing, token
counts, retrieval statistics, and stage-level attribution. Section~\ref{subsec:metrics} 
lists these fields and their definitions.

% % --- Block C: System Boundary ---

Each node runs a lightweight metrics module that reads hardware energy counters
at the start and end of every session. The orchestrator coordinates session
boundaries across all three nodes via a shared \texttt{session\_id}. Attribution
is at the device level (CPU package, GPU device); the framework does not
decompose energy to individual processes or kernels.
% Section~\ref{subsubsec:energy-metrics} defines the counter sources and
% computation method.

Excluded from the measurement boundary are power-at-the-wall measurements, network
fabric energy, storage I/O energy, and orchestrator computation. The orchestrator
performs no heavy computation; its role is control flow and data assembly.


% % --- Block D: Supporting Metrics ---
Beyond energy, each session records latency (time to first token, total
response time), token counts (query, augmentation, context, answer), retrieval
traces (pre-rerank, post-rerank, final context), and quality scores
(deterministic, provenance-based, and LLM-assisted). These supporting
measurements ensure that no energy claim is made without the corresponding
operational and quality context.


\subsection{Metrics Definition}
\label{subsec:metrics}
% This chapter defines what is measured and how. Interpretation belongs to
% Chapter~\ref{sec:baseline_analysis}. Section~\ref{subsec:metrics} specifies each metric's
% unit, computation method, and role in the analysis.

\subsubsection{Energy Metrics}
\label{subsubsec:energy-metrics}
% --- Block A: Primary energy quantities ---
Request energy represents the total device-level cost of answering one query
under one configuration. It is computed as the sum of the three infrastructure
nodes:
$$E_{\text{request}} = E_{\text{emb. model}} + E_{\text{database}} + E_{\text{llm}}$$

Each node contributes CPU and GPU energy:
$$E_{\text{node}} = E_{\text{cpu}} + E_{\text{gpu}}$$

where:
\begin{itemize}
    \item \textbf{$E_{cpu}$} is read from Intel RAPL counters, which report
          cumulative energy in microjoules ($\mu$J), converted to joules (J).
    \item \textbf{$E_{gpu}$} is read from the NVIDIA NVML interface, which
          reports cumulative energy in millijoules (mJ), converted to joules (J).
\end{itemize}

Both counter types are monotonic. Session energy is computed as a counter delta:
$\Delta E = E_{\text{end}} - E_{\text{start}}$ within the session window.

Energy is reported in joules rather than watts. Watts measure instantaneous
power draw; joules capture total energy consumed over the session window, which
matches the per-request comparison objective.

% --- Block B: Normalization (J/token) ---
In LLM systems, computation cost scales approximately with the number of tokens
processed. Because configurations change token counts per request, HyDE adds
augmentation tokens, reranking alters context length; normalizing energy by
total tokens controls for workload-size differences:

% writing_notes :
% allows comparison across: batch sizes, contexts, generation lengths, datasets
% Energy per token approximates:
% FLOPs per token × Energy per FLOP

$$E_{\text{token}} = \frac{E_{\text{request}}}{N_{\text{total\_tokens}}} \quad [\text{J/token}]$$

Joules per token serves as a comparison aid when output lengths differ across
configurations. It is not a perfect universal efficiency metric. It still ignores retrieval
cost amortization, query-level fixed energy overhead and short-answer bias but takes 


\subsubsection{Latency Metrics}
\label{subsubsec:performance-metrics}

% --- Block D: Latency ---
Two main latency measures capture the operational cost of each session:

\begin{figure}[h!]
    \centering
        \centering
        \includegraphics[width=0.8\textwidth]{latency.png}
        \caption{Latency measurement points}
        \label{fig:latency-measurement}
\end{figure}

\begin{itemize}
    \item \textbf{Time to first token (TTFT)} elapsed time in milliseconds
          from query submission to the first generated token. TTFT reflects
          user-perceived responsiveness and is sensitive to pre-generation stages
          such as HyDE and reranking.
    \item \textbf{Total response time} elapsed time in milliseconds from
          query submission to answer completion. This measures end-to-end duration.
\end{itemize}

Their difference can be defined as generation time. Per-stage timing intervals (embedding, 
retrieval, reranking, HyDE, generation)
are also recorded. They connect module activation to latency contributions and
support energy attribution by stage. 

% --- Block E: Token accounting ---

Each session records five token counts that describe workload size:
\begin{itemize}
    \item \textbf{query\_tokens}: tokens in the original user query
    \item \textbf{augm\_tokens}: tokens produced by HyDE (zero when HyDE
          disabled)
    \item \textbf{context\_tokens}: tokens in the retrieved context assembled
          for the LLM
    \item \textbf{answer\_tokens}: tokens in the generated answer
    \item \textbf{total\_tokens}: sum of the above four fields
\end{itemize}

Token counts serve two roles. First, they characterize how each configuration
changes workload composition (HyDE adds augmentation tokens; 
context length altered by reranking). 
Second, total tokens provides the denominator for energy
normalization (J/token) as defined earlier.


\subsubsection{Quality Proxies}
\label{subsubsec:quality-proxies}

The studied configurations activate modules that primarily affect retrieval
and context selection. Quality evaluation must therefore cover answer
correctness, retrieval correctness and grounding faithfulness.
No human evaluation is performed; all scores are computed by deterministic
text comparisons or LLM-assisted metrics and treated as proxies, not as
ground truth.

% --- HotpotQA benchmark metrics ---

\paragraph{HotpotQA benchmark metrics}~\cite{hotpotqa2018} provide
deterministic, reproducible scoring across the three evaluation axes:

\begin{itemize}
    \item \textbf{Answer~F1} token-level overlap between the generated
          answer and the reference answer after normalization (lowercasing,
          article and punctuation removal). Exact Match~(EM) is the
          corresponding binary variant.
    \item \textbf{Supporting Facts~F1} paragraph-level retrieval
          correctness: measures whether the pipeline retrieved the gold
          supporting-fact paragraphs annotated in the dataset.
    \item \textbf{Joint~F1} product of Answer~F1 and Supporting Facts~F1;
          rewards correct answers only when supported by correct retrieval.
\end{itemize}

Joint~F1 separates two failure modes: ``answer looks right but retrieval
failed'' versus ``answer is supported by the correct documents.'' This
distinction matters because the studied modules (HyDE, reranking, hybrid
search) primarily affect retrieval behavior.

% --- LLM-assisted quality (RAGAS) ---

\paragraph{LLM-assisted metrics} add semantic evaluation that string-based
metrics miss. The LLM-as-judge paradigm uses a language model to score
generated answers against references or contexts;
vRAG-Eval~\cite{vrageval2024} reports 82.6\% agreement between GPT-4
judgments and human annotations on RAG evaluation tasks. The thesis uses the
RAGAS framework~\cite{ragas2023} with two primary metrics, scored 0~to~1:

\begin{itemize}
    \item \textbf{Faithfulness} fraction of claims in the generated
          answer that are supported by the retrieved contexts.
    \item \textbf{Factual correctness} semantic agreement between
          the generated answer and the reference answer.
\end{itemize}

Faithfulness detects hallucination (unsupported claims); factual correctness
detects wrong answers (semantic mismatch with the reference). Together with
Answer~F1, they form the quality basis for the acceptability thresholds
defined in Chapter~\ref{sec:controller}. A dedicated evaluator model
(Qwen2.5-14B-Instruct), distinct from the pipeline LLM, scores
each answer; results are treated as assisted evaluation, not as
oracle truth.
% and scores should be interpreted comparatively across configurations (DO NOT DLT - keep commented)
% rather than as absolute quality measures. (DO NOT DLT - keep commented)


\subsection{Benchmark Dataset}
\label{subsec:dataset}
The measurement framework requires three aligned artifacts: 
a reference answer, gold supporting-fact annotations, and a retrieval corpus. 
This thesis utilizes HotpotQA~\cite{hotpotqa2018}, a multi-hop 
benchmark requiring information synthesis across multiple documents.

HotpotQA was selected for three primary reasons. First, its multi-hop 
structure mimics high-level knowledge management behavior, where users 
often seek to synthesize information across disparate internal sources. 
Second, the dataset is well-documented in existing literature with clear, 
comparable benchmarks, providing a rigorous baseline for our system's 
quality evaluation. Finally, multi-hop queries exercise the full RAG pipeline: 
retrieval must surface multiple documents and generation must synthesize across 
them. This complexity is a prerequisite for routing research, as it should create 
variation in module utility—some queries require advanced modules like 
reranking to succeed, while others do not.

While other knowledge-intensive datasets were considered, HotpotQA provided 
a better balance of reasoning depth and benchmarking clarity.

\begin{figure}[h!]
    \centering
    \includegraphics[width=0.85\textwidth]{hotpot.jpg}
    \caption{HotpotQA question and answer type distributions~\cite{hotpotqa2018}}
    \label{fig:hotpot}
\end{figure}

The full HotpotQA distractor split contains approximately 7,400 queries. 
However, executing each query under all eight configurations yields 59,200 
total sessions. To manage the computational demand and time constraints 
associated with running these sessions—specifically the high energy cost of 
LLM-assisted evaluation, we utilized a representative sample of 1,200 queries; 
Split equally between 600 bridge queries (requiring information 
from linked documents) and 600 comparison queries (comparing two entities). 
Each record contains the question text, a reference answer and gold 
supporting-fact paragraphs.

The document corpus consists of the specific Wikipedia pages referenced as 
gold supporting facts for these queries, which were chunked and pre-indexed in 
Milvus. This closed-corpus design ensures that every gold document is present 
in the retrieval index, allowing the Supporting Facts F1 metric to measure 
retrieval effectiveness directly rather than corpus coverage. The implications 
of this reduced retrieval difficulty on domain generalization are further addressed.

\subsection{Experimental Protocol}
\label{subsec:protocol}

% --- Block A: Protocol objective ---

Energy and latency measurements are sensitive to transient system state:
background load, thermal conditions, and caching effects can introduce
variation unrelated to configuration differences. The experimental protocol
controls for these factors by enforcing a fixed workload, a fixed
configuration per session, and deterministic scheduling. The same protocol is
reused when evaluating the controller in Chapter~\ref{sec:strategies}, so that
energy deltas between baseline and controlled execution remain comparable.

% --- Block B: Configuration space ---

The configuration space is the $2^3$ factorial design introduced in
Section~\ref{subsec:baseline-config}: three binary module axes (HyDE,
reranking, search type) yield eight configurations (C0--C7). Each
configuration is fixed for the duration of a session; there is no mid-request
switching. Table~\ref{tab:configs} lists all eight configurations with their
module settings.

\begin{table}[h!]
\centering
\begin{tabular}{c c c l l}
\hline
\textbf{Config} & \textbf{HyDE} & \textbf{Reranking} & \textbf{Search type} & \textbf{Description} \\
\hline
C0 & off & off & dense  & minimal pipeline \\
C1 & on  & off & dense  & + HyDE  \\
C2 & off & on  & dense  & + reranking \\
C3 & off & off & hybrid & + hybrid \\
C4 & on  & on  & dense  & + HyDE + reranking \\
C5 & on  & off & hybrid & + HyDE + hybrid \\
C6 & off & on  & hybrid & + reranking + hybrid \\
C7 & on  & on  & hybrid & all modules \\
\hline
\end{tabular}
\caption{Configuration space: $2^3$ factorial over three binary module axes.}
\label{tab:configs}
\end{table}

% --- Block C: Experimental unit and dataset size ---

The full experiment runs all 1{'}200~queries through all
8~configurations, producing 9{'}600~sessions. Each session logs raw energy
traces, timing metadata, token counts, retrieval traces and the generated
answer text. Runs order are shuffled to reduce systematic drift effects.
Post-hoc evaluation scripts then compute deterministic quality
metrics (Answer~F1, Supporting Facts~F1, Joint~F1) and LLM-assisted metrics
(RAGAS faithfulness, factual correctness) from the stored artifacts.

This dataset serves three purposes: benchmarking the CITI KMS system for
the first time against literature baselines, the baseline trade-off analysis
in Chapter~\ref{sec:baseline_analysis} and generate the training data for the
conditional controller in Chapters~\ref{sec:controller}--\ref{sec:strategies}.

% --- Block D: Execution scheduling ---

A checkpoint file records every completed session. If a run is interrupted, it
resumes from the last checkpoint without re-executing completed sessions. Failed
sessions are logged and excluded; they do not silently corrupt the dataset.

% --- Block E: Session procedure ---

Each session follows a fixed procedure:

\begin{enumerate}
    \item The orchestrator creates a unique \texttt{session\_id} and records the
          query identifier, configuration identifier and run number.
    \item The orchestrator broadcasts a \texttt{start} signal to all three
          infrastructure nodes. Each node begins reading energy counters and
          opens a session buffer.
    \item The pipeline executes the request under the fixed configuration:
          embedding, retrieval (dense or hybrid), optional HyDE, optional
          reranking, context assembly and generation.
    \item Stage marks (start and end timestamps) are emitted for each major
          pipeline stage. Marks are best-effort; a missed mark does not abort
          the session.
    \item After the answer is produced, the orchestrator broadcasts a
          \texttt{stop} signal. Each node closes its buffer and returns the
          collected traces.
    \item All artifacts: energy traces, session metadata, retrieval trace,
          answer text and stage intervals are persisted under the \texttt{session\_id}.
\end{enumerate}

Session energy is computed as a counter delta over the start-to-stop window.
% The framework does not subtract an idle baseline; because all configurations
% run under the same protocol, energy deltas between configurations remain
% meaningful without idle subtraction.

% --- Block F: Post-hoc evaluation ---

Evaluation is separated from execution. Deterministic metrics (Answer~F1,
Supporting Facts~F1, Joint~F1) are computed offline from the stored answer
texts and retrieval traces. RAGAS metrics (faithfulness, factual correctness)
are computed by a separate evaluation script that retrieves context documents
from Milvus and invokes a dedicated evaluator model (Qwen2.5-14B-Instruct),
distinct from the pipeline LLM to avoid self-evaluation bias and a perfomance ceiling. 

This separation keeps the online request path simple and allows the analysis to
be rerun. For example, with updated frameworks, metrics or corrected provenance
mappings, without re-executing all the sessions.

% Evaluation is separated from execution. Deterministic metrics (Answer~F1,
% Supporting Facts~F1, Joint~F1) are computed offline from the stored answer
% texts and retrieval traces. RAGAS metrics (faithfulness, factual correctness)
% are computed by a separate evaluation script that retrieves context documents
% from Milvus.

% This separation keeps the online request path simple and allows the analysis to
% be rerun. For example, with updated frameworks, metrics or corrected provenance
% mappings, without re-executing all the sessions.

% Also it invokes a dedicated evaluator model (Qwen2.5-14B-Instruct),
% distinct from the pipeline LLM to avoid self-evaluation bias.
% Recent benchmarks show that evaluator reliability scales with model size and
% that instruction-tuned models produce more consistent judgments on open-ended
% quality metrics~\cite{llmasjudge2024}; we therefore select a 14B instruct model
% as the largest evaluator our infrastructure can serve alongside the pipeline.

% --- Block G: Validity criteria ---

Several validity checks are applied before any session enters the analysis
dataset. Energy counters must be monotonic within a session; a negative delta
invalidates that node's energy for the session. Each session must contain
start and stop boundaries for every node included in the energy total. Sessions
that fail any check are flagged and excluded from the analysis. No silent
corrections are applied; exclusion is preferred over ad-hoc imputation.

% --- Block H: Protocol invariance ---

All experiments in this thesis: baseline analysis controller training
and controller evaluation use the same query set, the same corpus and
index, the same eight configurations, the same measurement boundary and the
same logging artifacts.


\subsection{Features}
\label{subsec:router-features}

% --- Block A: Feature record policy ---

The conditional controller proposed predicts a
configuration \emph{before} the pipeline executes. This constraint divides the
logged data into two categories:

\begin{itemize}
    \item \textbf{Pre-execution features} are derived from the query text alone.
          They include :
          \begin{itemize}
            \item \textbf{surface statistics} character and token counts, sentence count)
            \item \textbf{question-type indicators} question word, definition or list intent, logical operators
            \item \textbf{complexity signals} named-entity count and diversity, noun-phrase count, content-word ratio
            \item \textbf{semantic features} embedding norm, cosine similarities to query-type centroids
          \end{itemize}
          These 23~features are computable before any pipeline stage runs.
    \item \textbf{Post-execution features} are produced by the pipeline and these are available only after the session completes.
    \begin{itemize}
            \item \textbf{retrieval scores } top-1 similarity, score gaps
            \item \textbf{context statistics} context length, number of documents, truncation ratio
            \item \textbf{token counts} 
            \item \textbf{quality metrics} 
      \end{itemize}

\end{itemize}

Pre-execution features serve as router inputs. Post-execution features are
logged for analysis but cannot serve as controller inputs without running
(and paying the cost of) at least part of the pipeline.
Chapter~\ref{sec:controller} formalizes this distinction and defines the
feature groups used for training.

% --- Block B: Feature persistence ---

Feature extraction is performed offline: a dedicated script computes all
23~pre-execution features for the 1{'}200~queries and stores them alongside the
query identifiers. The same session artifacts that feed the baseline analysis
also feed the controller training pipeline. This avoids a second, inconsistent
data-collection step and ensures that the router is trained on the exact system
behavior measured in this chapter.








% % ============================================================================
% DO NOT DELETE  : 
% Factoid queries
% Straightforward, single-hop questions that can typically be answered from one document / one passage.

% Multi-document queries
% Questions that intentionally require retrieving information from multiple documents and combining them (e.g., compare, aggregate, reconcile).

% Reasoning / implication queries
% Queries where the answer requires inference beyond direct lookup, using retrieved facts to draw a conclusion, explain consequences, or propose a next step.

% Ambiguous / underspecified queries
% Real users ask: “How does this work again?” or “What’s the process?”
% These are where HyDE/query expansion often matters.

% Negative / no-answer queries
% Questions that shouldn’t be answered from the corpus (not present, out of scope).
% This is crucial if your controller also decides “don’t retrieve” or “retrieve but abstain”.

% If you add just (2) underspecified and (3) no-answer, your dataset becomes much more “future-proof” for real traffic.

% A sharper way to frame the dataset

% Instead of “3 types of questions”, you can describe it as training coverage across two axes:

% Retrieval breadth: single-doc → multi-doc

% Answer processing: extractive → compositional → inferential

% Your three buckets are basically corners of this grid—good! The improvement is making sure you populate the in-between regions so the classifier doesn’t overfit.

% ---

% We evaluate our retrieval-augmented generation system using the KILT benchmark (Petroni et al., NAACL 2021), which provides a standardized Wikipedia corpus, curated query sets, and gold provenance annotations. This allows rigorous evaluation of both retrieval correctness and answer quality while enabling controlled system-level measurements.
% -->  if we Use KILT as the academic reference Add your own domain corpus (lab documents)
%     You compare: Δ system behavior under identical queries. This is key.
%             “We evaluate our RAG system on KILT, a standardized benchmark for knowledge-intensive NLP tasks grounded in Wikipedia.”

%             You are no longer benchmarking QA performance

% You are benchmarking system behavior under retrieval uncertainty.

% That’s a different but valid scientific claim.

% Aspect	kilt_wikipedia	kilt_corpus
% QA benchmark fidelity	High	Lower
% Retrieval difficulty	Artificially easier	Realistic
% Noise	Low	High
% Router value	Moderate	High
% Energy signal	Muted	Strong
% Systems relevance	Medium	Very high

% With kilt_corpus, your router actually matters more.

% --> “While KILT provides a task-aligned Wikipedia subset, we intentionally use the full KILT corpus to reflect realistic retrieval conditions. Our objective is not to maximize QA benchmark scores, but to study energy–latency–quality trade-offs in retrieval-augmented systems under large-scale retrieval noise.”



% use Kepler (Kubernetes-based Efficient Power Level Exporter) to collect energy consumption data 2601.02522v2.pdf